% Figure: charged-particle motion in a magnetic field. Left: pure gyration
% (B out of page), opposite senses for ion and electron. Right: cycloidal
% E x B drift, the guiding centre marching steadily perpendicular to both
% E and B with the same velocity for both signs of charge.
\begin{figure}[htbp]
\centering
\begin{tikzpicture}[>=latex, scale=1.0]
% ===== Left panel: gyration =====
\begin{scope}[shift={(0,0)}]
  \node[font=\small] at (0,2.6) {gyration ($\vect{B}$ out of page)};
  % B field dots
  \foreach \x in {-1.6,-0.8,0,0.8,1.6}{
    \foreach \y in {-1.6,-0.8,0.8,1.6}{
      \node[engblue, font=\tiny] at (\x,\y) {$\odot$};
    }
  }
  % ion orbit (large, clockwise)
  \draw[engred, thick] (0.9,0) circle (0.9);
  \draw[engred, ->, thick] (0.9,0.9) -- (1.0,0.9);
  \node[engred, font=\scriptsize] at (0.9,0) {ion};
  % electron orbit (small, counter-clockwise)
  \draw[englime!60!black, thick] (-1.0,0) circle (0.45);
  \draw[englime!60!black, ->, thick] (-1.0,0.45) -- (-1.1,0.45);
  \node[englime!60!black, font=\scriptsize] at (-1.0,-0.85) {e$^-$};
\end{scope}
% ===== Right panel: E x B drift =====
\begin{scope}[shift={(6.2,0)}]
  \node[font=\small] at (0,2.6) {$\vect{E}\times\vect{B}$ drift};
  % E field arrows pointing up
  \draw[engorange, ->, thick] (-2.2,-1.8) -- (-2.2,1.8);
  \node[engorange, font=\scriptsize] at (-2.5,1.6) {$\vect{E}$};
  \node[engblue, font=\tiny] at (-1.6,-1.8) {$\odot\,\vect{B}$};
  % cycloid (guiding centre drifts to the right)
  \draw[engblue, thick, samples=160, domain=0:9, variable=\t]
    plot ({-1.4 + 0.42*\t - 0.42*sin(deg(\t))}, {0.42 - 0.42*cos(deg(\t))});
  % drift arrow
  \draw[enggray, ->, thick] (-1.4,-1.4) -- (1.9,-1.4)
    node[midway, below, font=\scriptsize] {$v_{E\times B}=E/B$};
\end{scope}
\end{tikzpicture}
\caption[Gyration and $E\times B$ drift]{Single-particle motion in a
magnetic field. Left: with $\vect{B}$ out of the page, ions and electrons
gyrate in opposite senses; the electron orbit is far smaller because the
Larmor radius scales with mass. Right: an added perpendicular electric
field bends each gyro-orbit into a cycloid whose guiding centre drifts at
$v_{E\times B}=E/B$, perpendicular to both $\vect{E}$ and $\vect{B}$ and
independent of charge, mass, and energy, so ions and electrons drift
together with no net current.}
\label{fig:vol04ch13:gyration-drift}
\end{figure}
